\chapter{曲线积分、曲面积分与场论初步}\label{ch:8}
\begin{tcolorbox}[title=章首知识地图,colframe=Ink]
弧长与面积度量 $\to$ 切向做功与法向通量 $\to$ Green、Gauss、Stokes $\to$ 势函数、散度和旋度。
\end{tcolorbox}
\begin{examtip}考研主线：两类曲线与曲面积分、Green、Gauss、Stokes及路径无关。常见题型：直接法与公式法、补线补面、方向判断。竞赛拓展：奇点与区域孔洞、势函数、散度旋度恒等式。上述为学习任务分层，具体范围以所报考科目的当年大纲为准。\end{examtip}
\section{两类曲线积分与参数化}\label{sec:8:1}
\subsection{概念、定理与推导}

\begin{definition}{弧长积分与做功积分}第一类曲线积分$\int_Lf\dd s$按长度加权，与方向无关；第二类$\int_L\bm F\cdot\dd\bm r=\int_LP\dd x+Q\dd y+R\dd z$依赖方向。\end{definition}
\begin{theorem}{参数计算}设$\bm r(t)$在$[a,b]$上分段$C^1$，$f$和$\bm F$在曲线上连续。沿参数描绘的曲线积分为
\begin{align}\int_Lf\dd s&=\int_a^bf(\bm r(t))|\bm r\prime(t)|\dd t,\\
\int_L\bm F\cdot\dd\bm r&=\int_a^b\bm F(\bm r(t))\cdot\bm r\prime(t)\dd t.\end{align}
\end{theorem}
\begin{proof}小线段长度近似为$|\bm r\prime|\Delta t$，有向位移近似为$\bm r\prime\Delta t$，分别代入积分和得到。\end{proof}
两类关系为$\bm F\cdot\dd\bm r=(\bm F\cdot\bm T)\dd s$，其中$\bm T$是方向单位切向量。

\begin{method}{两类曲线积分与参数化}条件：参数方向与题设一致。步骤：代入坐标，计算速度模或速度向量；分段曲线分段积分。易错点：第一类用模长，第二类保留符号。\end{method}
\subsection{例题与方法比较}
\begin{example}{基础例题｜圆周弧长}{难度：$\star$}
\label{ex:8:1}\textbf{题目} 单位圆上求$\int_Lx^2\dd s$。\par
\textbf{思路} 参数化或对称。\par
\textbf{解答} 方法一：$x=\cos t,y=\sin t,\dd s=\dd t$，积分$\int_0^{2\pi}\cos^2t\dd t=\pi$。方法二：由$x,y$对称两平方积分相等，其和为周长$2\pi$，每项$\pi$。\par
\textbf{方法点评} 对称法避免三角积分。\par
\textbf{易错点或变式} 反向走圆周第一类积分仍为$\pi$。
\end{example}
\begin{example}{考研风格例题｜圆周做功}{难度：$\star\star$}
\label{ex:8:2}\textbf{题目} 逆时针单位圆求$\oint_L(-y\dd x+x\dd y)$。\par
\textbf{思路} 直接参数化或Green。\par
\textbf{解答} 方法一：代入后被积式$(\sin^2t+\cos^2t)\dd t$，值$2\pi$。方法二：$Q_x-P_y=2$，Green给$2\iint_D1\dd A=2\pi$。\par
\textbf{方法点评} 直接法适合规则曲线，Green适合闭曲线与简单旋度。\par
\textbf{易错点或变式} 顺时针结果为$-2\pi$。
\end{example}
\begin{example}{考研风格例题｜保守场}{难度：$\star\star$}
\label{ex:8:3}\textbf{题目} 任意曲线从$(0,0)$至$(1,2)$，求$\int_L(2x\dd x+2y\dd y)$。\par
\textbf{思路} 识别全微分。\par
\textbf{解答} 方法一：势函数$\varphi=x^2+y^2$，积分为$5-0=5$。方法二：选折线先沿$x$到$(1,0)$再沿$y$到$(1,2)$，两段为1与4，总和5；路径无关由势函数保证。\par
\textbf{方法点评} 势函数法给任意路径结论，选路径只在已知路径无关后合法。\par
\textbf{易错点或变式} 不能未经证明随意把题设曲线换成折线。
\end{example}
\section{Green公式与路径无关}\label{sec:8:2}

\begin{center}\begin{tikzpicture}[scale=.8,>=stealth]
\fill[Teal!10](0,0)circle(1.6);\fill[white](0,0)circle(.65);
\draw[Ink,thick](0,0)circle(1.6);\draw[Ink,thick](0,0)circle(.65);
\draw[->,Teal,very thick](1.6,0)arc(0:60:1.6);
\draw[->,Gold,very thick](.65,0)arc(0:-110:.65);
\node[right]at(2,1){外边界逆时针};\node[right]at(2,-.5){内边界顺时针};
\end{tikzpicture}\par\small 两条边界均保持区域在行进方向左侧。
\end{center}

\subsection{概念、定理与推导}

\begin{theorem}{Green公式}设$D$为边界分段光滑的有界平面区域，$P,Q$在包含闭包$\overline D$的一个开邻域内为$C^1$，则
\begin{equation}\oint_{\partial D}P\dd x+Q\dd y=\iint_D(Q_x-P_y)\dd A.\end{equation}
正向要求行进时区域在左；有孔时外圈逆时针、内圈顺时针。\end{theorem}
\begin{proof}对可投影简单区域分别把$Q_x$和$-P_y$积分化成两侧边界差；一般区域分块，内部公共边方向相反而抵消。\end{proof}
\begin{theorem}{路径无关的常用判据}在开连通区域，连续场有势函数等价于分段光滑曲线积分只依赖端点；对于$C^1$平面场，在单连通区域$P_y=Q_x$可保证势函数存在。\end{theorem}
\begin{proof}有势时沿曲线使用链式法则。反向固定基点用路径积分定义势，路径无关使定义唯一。单连通与Green使所有闭路积分为零。\end{proof}

\begin{method}{Green公式与路径无关}条件：场在区域及内部无奇点。步骤：确认闭合与方向，算$Q_x-P_y$；非闭曲线可补线并减去补线积分。易错点：旋度零但区域有孔时不保证路径无关。\end{method}
\subsection{例题与方法比较}
\begin{example}{考研风格例题｜面积公式}{难度：$\star\star$}
\label{ex:8:4}\textbf{题目} 求椭圆逆时针边界上的$\oint(x\dd y-y\dd x)$，半轴$a,b>0$。\par
\textbf{思路} Green或椭圆参数。\par
\textbf{解答} 方法一：Green给$2\operatorname{Area}(D)=2\pi ab$。方法二：$x=a\cos t,y=b\sin t$，表达式恒为$ab\dd t$，积分$2\pi ab$。\par
\textbf{方法点评} 面积公式从常旋度产生。\par
\textbf{易错点或变式} 若题目是面积，须再除以2。
\end{example}
\begin{example}{考研风格例题｜有孔反例}{难度：$\star\star$}
\label{ex:8:5}\textbf{题目} 平面去掉原点，$\bm F=(-\frac{y}{r^2},\frac{x}{r^2})$，判断是否全局有势。\par
\textbf{思路} 局部旋度零但圆周积分非零。\par
\textbf{解答} 方法一：单位圆上积分$\int_0^{2\pi}1\dd t=2\pi$，所以不存在全局单值势。方法二：局部势为极角$\theta$，绕原点一周增加$2\pi$，不能在穿孔平面定义连续单值角函数。\par
\textbf{方法点评} 拓扑障碍不是求偏导能发现的。\par
\textbf{易错点或变式} 不能对含奇点原点的圆盘使用Green。
\end{example}
\begin{example}{考研风格例题｜补线}{难度：$\star\star$}
\label{ex:8:6}\textbf{题目} 沿上半单位圆从$(1,0)$到$(-1,0)$求$\int_L y\dd x$。\par
\textbf{思路} 加$x$轴线段闭合。\par
\textbf{解答} 补线从$(-1,0)$到$(1,0)$，其积分为0，整体逆时针。$P=y,Q=0$，$Q_x-P_y=-1$，积分为负半圆面积$-\frac{\pi}{2}$。直接参数$t\in[0,\pi]$给$-\int_0^\pi\sin^2t\dd t$同值。\par
\textbf{方法点评} 补线选让积分易算的边界。\par
\textbf{易错点或变式} 判断闭合后的方向，不能只看弧的形状。
\end{example}
\section{两类曲面积分与通量}\label{sec:8:3}

\begin{method}{三种投影的定向面积元}
若$z=g(x,y)$，朝上取$(-g_x,-g_y,1)\dd x\dd y$；若$x=g(y,z)$，朝正$x$方向取$(1,-g_y,-g_z)\dd y\dd z$；若$y=g(z,x)$，朝正$y$方向取$(-g_x,1,-g_z)\dd z\dd x$。反向全部乘$-1$。这些向量已包含面积倍率，不再额外乘$\sqrt{1+|\nabla g|^2}$。第一类积分则只取对应向量的模长。
\end{method}

\subsection{概念、定理与推导}

\begin{definition}{曲面积分}第一类$\iint_Sf\dd S$对面积加权，与定向无关。第二类通量$\iint_S\bm F\cdot\bm n\dd S$依赖指定单位法向；反向变号。记作$\iint_SP\dd y\dd z+Q\dd z\dd x+R\dd x\dd y$时三个投影面积元带方向。\end{definition}
\begin{theorem}{参数与图形公式}参数曲面$\bm r(u,v)$满足$\bm r_u\times\bm r_v\ne0$时，$\dd S=|\bm r_u\times\bm r_v|\dd u\dd v$。对$z=g(x,y)$，上侧有$\bm n\dd S=(-g_x,-g_y,1)\dd x\dd y$，$\dd S=\sqrt{1+g_x^2+g_y^2}\dd x\dd y$。\end{theorem}
\begin{proof}小参数矩形被切向量映成小平行四边形，叉积模长给面积，叉积向量给定向面积元。\end{proof}

\begin{method}{两类曲面积分与通量}条件：先确认曲面投影与方向。步骤：第一类用叉积模，第二类用定向叉积；图形曲面优先投影。易错点：通量公式再乘一次面积根号会重复计量。\end{method}
\subsection{例题与方法比较}
\begin{example}{基础例题｜平面片面积}{难度：$\star$}
\label{ex:8:7}\textbf{题目} 求$z=x+y$在$0\le x,y\le1$上曲面面积。\par
\textbf{思路} 图形面积公式。\par
\textbf{解答} $g_x=g_y=1$，面积$\int_0^1\int_0^1\sqrt3\dd x\dd y=\sqrt3$。参数边向量$(1,0,1),(0,1,1)$叉积$(-1,-1,1)$模同为$\sqrt3$。\par
\textbf{方法点评} 线性参数面是平行四边形，几何与积分一致。\par
\textbf{易错点或变式} 投影面积1与实际面积$\sqrt3$不同。
\end{example}
\begin{example}{考研风格例题｜平面通量}{难度：$\star\star$}
\label{ex:8:8}\textbf{题目} $\bm F=(0,0,1)$穿过上一题曲面上侧的通量。\par
\textbf{思路} 上侧定向面积元。\par
\textbf{解答} $\bm F\cdot(-1,-1,1)=1$，积分为1。也可单位法向的$z$分量为$\frac{1}{\sqrt3}$，乘曲面面积$\sqrt3$得到1。\par
\textbf{方法点评} 恒定竖直场通量等于投影面积。\par
\textbf{易错点或变式} 下侧通量为$-1$。
\end{example}
\begin{example}{考研风格例题｜球面径向通量}{难度：$\star\star$}
\label{ex:8:9}\textbf{题目} 半径$R$球面外侧，求$\bm F=(x,y,z)$的通量。\par
\textbf{思路} 场与单位外法向同向。\par
\textbf{解答} $\bm n=\frac{\bm r}{R}$，$\bm F\cdot\bm n=R$，通量$4\pi R^3$。Gauss法散度为3，体积$\frac{4\pi R^3}{3}$，同得。\par
\textbf{方法点评} 直接法利用恒定法向分量，公式法避免参数化。\par
\textbf{易错点或变式} 通量的量纲是长度三次，不是球面积。
\end{example}
\section{Gauss公式、散度与补面}\label{sec:8:4}
\subsection{概念、定理与推导}

\begin{definition}{散度}对$\bm F=(P,Q,R)$，$\diver\bm F=P_x+Q_y+R_z$，描述单位体积净向外通量的局部密度。\end{definition}
\begin{theorem}{Gauss散度定理}对分段光滑闭曲面包围的有界区域$\Omega$，若$\bm F$在包含$\overline\Omega$的一个开邻域内为$C^1$，则
\begin{equation}\iint_{\partial\Omega}\bm F\cdot\bm n\dd S=\iiint_\Omega\diver\bm F\dd V,\end{equation}
其中$\bm n$为外法向。\end{theorem}
\begin{proof}先对小长方体，将相对两面的通量差用一维基本定理变为偏导体积分；拼合时内部面通量抵消，取极限得一般区域结论。\end{proof}
闭曲面是必要几何条件；开曲面可补面，求闭合通量后扣除补面。场的奇点不能置于所用体区域内而忽略。

\begin{method}{Gauss公式、散度与补面}条件：闭合、外向、无内部奇点。步骤：算散度，选易积分的体域；开面加底面后扣除。易错点：补面的外法向可能向下。\end{method}
\subsection{例题与方法比较}
\begin{example}{考研风格例题｜半球补面}{难度：$\star\star$}
\label{ex:8:10}\textbf{题目} 上半单位球面取外侧，求$\bm F=(x,y,z)$的通量。\par
\textbf{思路} 补$z=0$底盘。\par
\textbf{解答} 闭区域体积$\frac{2\pi}{3}$，散度3，总通量$2\pi$。底盘外法向$(0,0,-1)$，场的法向分量$-z=0$，故半球通量$2\pi$。直接法球面上法向分量1，半球面积$2\pi$，同得。\par
\textbf{方法点评} 补面先算其贡献，不能凭图形假定为零。\par
\textbf{易错点或变式} 底盘为零依赖具体场。
\end{example}
\begin{example}{考研风格例题｜立方体}{难度：$\star\star$}
\label{ex:8:11}\textbf{题目} 单位立方体$[0,1]^3$外侧求$\bm F=(x^2,y^2,z^2)$通量。\par
\textbf{思路} 散度法或逐面法。\par
\textbf{解答} 散度$2(x+y+z)$，体积分为3。逐面法：三个坐标为0的面通量0，三个坐标为1的面各通量1，总3。\par
\textbf{方法点评} 多面体场简单时逐面法便于核验。\par
\textbf{易错点或变式} 负向外法向对应低坐标面。
\end{example}
\begin{example}{考研风格例题｜源点场}{难度：$\star\star$}
\label{ex:8:12}\textbf{题目} 半径$R$球面外侧求$\bm F=\frac{\bm r}{|\bm r|^3}$通量。\par
\textbf{思路} 直接算，避免在含原点球体上误用Gauss。\par
\textbf{解答} 球面上$\bm F\cdot\bm n=\frac{1}{R^2}$，乘面积$4\pi R^2$得$4\pi$。虽然原点外散度为0，但场在原点未定义；对球壳应用Gauss只说明内外球面的向外径向通量相等。\par
\textbf{方法点评} 奇点决定非零总通量。\par
\textbf{易错点或变式} 不能以“散度为0”直接得闭球通量0。
\end{example}
\section{Stokes、旋度与场论联系}\label{sec:8:5}
\subsection{概念、定理与推导}

\begin{definition}{旋度}$\rot\bm F=(R_y-Q_z,P_z-R_x,Q_x-P_y)$，表示局部环流的方向与强度。\end{definition}
\begin{theorem}{Stokes公式}若$S$是带分段光滑边界的紧致定向分段光滑曲面，$\bm F$在包含$S$的一个开邻域内为$C^1$，边界方向与法向满足右手规则，则
\begin{equation}\oint_{\partial S}\bm F\cdot\dd\bm r=\iint_S(\rot\bm F)\cdot\bm n\dd S.\end{equation}
\end{theorem}
\begin{proof}将曲面参数化，在参数域对切向一形式用Green；链式法则整理后右端为旋度与定向叉积的内积。多片拼合时内部边界抵消。\end{proof}
\begin{extension}对二阶连续可微函数，$\rot(\nabla f)=0$；对二阶连续向量场，$\diver(\rot\bm F)=0$。两式均由混合偏导交换而抵消。全局“无旋即有势”仍需区域拓扑条件。\end{extension}

\begin{method}{Stokes、旋度与场论联系}条件：边界与法向右手一致，场在所选曲面邻域光滑。步骤：闭空间曲线可换成同边界的简单曲面；计算旋度投影。易错点：随意换面跨过场的奇点。\end{method}
\subsection{例题与方法比较}
\begin{example}{竞赛风格例题｜空间圆环流}{难度：$\star\star\star$}
\label{ex:8:13}\textbf{题目} 曲线$x^2+y^2=1,z=2$从上看逆时针，求$\oint(-y\dd x+x\dd y+z\dd z)$。\par
\textbf{思路} 选水平圆盘。\par
\textbf{解答} 旋度$(0,0,2)$，上法向对应题设方向，积分$2\pi$。直接参数$(\cos t,\sin t,2)$得被积式$\dd t$，同为$2\pi$。\par
\textbf{方法点评} 曲线高度对该环流无影响，$z\dd z$为全微分。\par
\textbf{易错点或变式} 从下看同一行进方向为顺时针。
\end{example}
\begin{example}{竞赛风格例题｜势函数重建}{难度：$\star\star\star$}
\label{ex:8:14}\textbf{题目} $\bm F=(yz,xz,xy)$，求从$(1,1,1)$到$(2,3,4)$的曲线积分。\par
\textbf{思路} 识别乘积梯度。\par
\textbf{解答} $\bm F=\nabla(xyz)$，故积分$24-1=23$。若逐坐标折线计算，三段分别1、4、18，合计23。\par
\textbf{方法点评} 势函数把曲线计算压缩成端点代入。\par
\textbf{易错点或变式} 无须知道实际路径形状，因为场在全空间有势。
\end{example}
\begin{example}{竞赛风格例题｜旋度通量}{难度：$\star\star\star$}
\label{ex:8:15}\textbf{题目} 设$\bm F\in C^2$，证明任意适用Gauss的闭曲面上$\iint\rot\bm F\cdot\bm n\dd S=0$。\par
\textbf{思路} 散度作用于旋度恒为零。\par
\textbf{解答} 由Gauss，左边等于$\iiint\diver(\rot\bm F)\dd V$。展开得到$R_{yx}-Q_{zx}+P_{zy}-R_{xy}+Q_{xz}-P_{yz}=0$，混合偏导成对抵消，故积分0。\par
\textbf{方法点评} 场论恒等式是微分运算与边界积分的兼容关系。\par
\textbf{易错点或变式} $C^2$保证此处混合偏导可交换。
\end{example}
\begin{warning}第一类积分不随方向改变，第二类积分反向变号；Green与Gauss的内部光滑条件不能省略；无旋的局部结论不能自动提升为全局势函数结论。\end{warning}


\section{场论恒等式、势与拓扑}\label{mat:fields}
\begin{examtip}将散度旋度、局部与全局势、通量源点、保守场进一步整理为可计算内容。\end{examtip}
\subsection{构造势函数的逐变量积分}\label{mat:fields:potential}
\begin{method}{构造势函数的逐变量积分}\textbf{适用条件} 全空间或适当单连通区域内无旋场。\par\textbf{实施与依据} 先积分一个分量，余函数依次由其他分量确定。\end{method}
\begin{example}{专题例题｜构造势函数的逐变量积分}{难度：$\star\star$}
\textbf{题目} 求$\bm F=(2xy+z,x^2, x)$的势。\par\textbf{思路与解答} 由$\varphi_x=2xy+z$得$\varphi=x^2y+xz+h(y,z)$。第二分量给$h_y=0$，第三给$h_z=0$，故$\varphi=x^2y+xz+C$。\par\textbf{方法点评与边界} 先确认交叉偏导相容，且势的全局存在依赖区域。
\end{example}
\subsection{Laplacian与径向场}\label{mat:fields:laplacian}
\begin{method}{Laplacian与径向场}\textbf{适用条件} 二阶可微标量函数或径向势。\par\textbf{实施与依据} $\Delta f=f_{xx}+f_{yy}+f_{zz}$；对$f=u(r)$，三维$\Delta f=u\prime\prime+\frac{2u\prime}{r}$，$r>0$。\end{method}
\begin{example}{专题例题｜Laplacian与径向场}{难度：$\star\star$}
\textbf{题目} 验证$\frac{1}{r}$在原点外调和。\par\textbf{思路与解答} $u\prime=-r^{-2},u\prime\prime=2r^{-3}$，合并为0。其梯度为$-\frac{\bm r}{r^3}$，原点是奇点，不能忽略通量源。\par\textbf{方法点评与边界} 原点外调和不等于在整个球内二阶可微。
\end{example}
\subsection{向量恒等式与边界形式}\label{mat:fields:curlgrad}
\begin{method}{向量恒等式与边界形式}\textbf{适用条件} 二阶连续偏导。\par\textbf{实施与依据} $\rot\grad f=0$、$\diver\rot\bm F=0$由混合偏导抵消；$\rot\rot\bm F=\grad\diver\bm F-\Delta\bm F$逐分量验证。\end{method}
\begin{example}{专题例题｜向量恒等式与边界形式}{难度：$\star\star$}
\textbf{题目} 若$\bm F=\grad f$，说明闭路做功。\par\textbf{思路与解答} 链式法则给沿任意闭路$\oint\grad f\cdot\dd\bm r=f(\text{终点})-f(\text{起点})=0$；这比仅计算旋度为零更直接地证明全局路径无关。\par\textbf{方法点评与边界} “无旋”反向推出“有势”仍需拓扑条件。
\end{example}




\section{曲线曲面积分与场论的直接法核对}\label{deep:fields_source_deep}

\begin{examtip}每次使用Green、Gauss或Stokes前先检查方向、边界与奇点。\end{examtip}

\subsection{两类曲线积分、Green与路径无关}\label{deep:fields_source_deep:green}

\begin{method}{两类曲线积分、Green与路径无关}\textbf{适用条件} 第一类按弧长，第二类按有向微分；Green要求平面区域内场正则。\par\textbf{步骤与依据} 直接参数化作核对；封闭路径优先看旋度；有洞区域补边界或检验环路积分。\par\textbf{易错点} 偏导相等仅给局部闭性，单连通或环路条件决定全局势函数。\end{method}

\begin{example}{考研风格｜圆周积分的直接与Green法}{难度：$\star\star$}\label{deep:ex:fields_source_deep:greenarea}\textbf{题目} $C:x^2+y^2=R^2$逆时针，求$\oint_C(-y\dd x+x\dd y)$。\par\textbf{思路} 参数圆周，或使用Green面积公式。\par\textbf{详细解答} 直接法$x=R\cos t,y=R\sin t$，被积式$R^2(\sin^2t+\cos^2t)\dd t$，积分$2\pi R^2$。Green法$Q_x-P_y=1-(-1)=2$，故为$2\iint_D\dd A=2\pi R^2$。\par\textbf{方法点评} 同一题同时验证方向与系数2。\par\textbf{变式与易错点} 若用$\frac12(x\dd y-y\dd x)$，结果才是面积。\end{example}

\begin{example}{竞赛风格｜穿孔区域中的闭而非恰当}{难度：$\star\star\star$}\label{deep:ex:fields_source_deep:hole}\textbf{题目} 在$\R^2\setminus\{0\}$考察$\omega=-\frac{y\dd x}{x^2+y^2}+\frac{x\dd y}{x^2+y^2}$。\par\textbf{思路} 偏导相等但绕原点一周积分非零。\par\textbf{详细解答} 直接算$P_y=Q_x=\frac{y^2-x^2}{(x^2+y^2)^2}$，所以旋度0。沿单位圆$\omega=\dd t$，环路积分$2\pi$，故不存在全域单值势函数。在不穿过一条射线的单连通子域，极角$\theta$可连续选支，$\omega=\dd\theta$。\par\textbf{方法点评} 展示拓扑条件而非代数计算决定路径无关。\par\textbf{变式与易错点} 不能对包含原点的圆盘直接套Green，因为场在原点未定义。\end{example}

\subsection{通量、补面与奇点}\label{deep:fields_source_deep:gauss}

\begin{method}{通量、补面与奇点}\textbf{适用条件} Gauss要求闭曲面和内部场有连续偏导；开放曲面可补简单面。\par\textbf{步骤与依据} 先散度；补面法写清所求通量等于闭面通量减补面通量；有奇点挖小球。\par\textbf{易错点} 外法向与题设上侧、下侧不一定一致。\end{method}

\begin{example}{考研风格｜半球通量双解}{难度：$\star\star\star$}\label{deep:ex:fields_source_deep:hemisphere}\textbf{题目} $S$为半径$R$上半球外侧，$F=(x,y,z)$，求通量。\par\textbf{思路} 补底盘用Gauss，或球面法向直接点乘。\par\textbf{详细解答} 补$z=0$底盘，其向外法向向下，$F\cdot n=-z=0$。散度3，半球体积$\frac{2\pi R^3}{3}$，所以通量$2\pi R^3$。直接法球面$n=\frac{1}{R}(x,y,z)$，$F\cdot n=R$，乘上半球面积$2\pi R^2$也得同值。\par\textbf{方法点评} 对径向场，直接法甚至比补面法更短。\par\textbf{变式与易错点} 若场的第三分量有常数项，底盘通量一般不为0。\end{example}

\begin{example}{竞赛风格｜反平方场的奇点通量}{难度：$\star\star\star$}\label{deep:ex:fields_source_deep:inverse}\textbf{题目} 求$F=\frac{\bm r}{|\bm r|^3}$通过任意包围原点的光滑闭曲面的外通量。\par\textbf{思路} 挖去原点小球，在剩余区域散度为0。\par\textbf{详细解答} 对$r\ne0$计算$\operatorname{div}F=3r^{-3}-3r^{-5}(x^2+y^2+z^2)=0$。挖半径$\varepsilon$球，穿孔域内Gauss给外曲面通量等于小球按其自身外法向的通量。球上$F\cdot n=\varepsilon^{-2}$，面积$4\pi\varepsilon^2$，故答案$4\pi$。若曲面不包原点且区域无奇点，则通量0。\par\textbf{方法点评} 公式法与球面对称直接法结合，得到形状无关性。\par\textbf{变式与易错点} 散度为0只在穿孔空间成立，不能忽略原点直接得总通量0。\end{example}

\subsection{Stokes、旋度与曲面替换}\label{deep:fields_source_deep:stokes}

\begin{method}{Stokes、旋度与曲面替换}\textbf{适用条件} 边界曲线方向与法向满足右手规则，场在曲面邻域正则。\par\textbf{步骤与依据} 先算旋度；在相同有向边界下选择最简单曲面；直接参数化边界作复核。\par\textbf{易错点} 曲面可以替换，但边界方向必须保留。\end{method}

\begin{example}{考研风格｜空间环路积分三解接口}{难度：$\star\star\star$}\label{deep:ex:fields_source_deep:stokes}\textbf{题目} $C$为$z=1,x^2+y^2=1$从上看逆时针，求$\oint_C(-y\dd x+x\dd y+z\dd z)$。\par\textbf{思路} 直接参数、Stokes平盘、全微分消项。\par\textbf{详细解答} 直接法$z=1,\dd z=0$，圆周给$2\pi$。Stokes法$F=(-y,x,z)$，$\operatorname{rot}F=(0,0,2)$，上向单位圆盘面积$\pi$，通量$2\pi$。第三观察$z\dd z=\dd(\frac{z^2}{2})$沿任意闭路积分0，问题退回平面Green公式。\par\textbf{方法点评} 对同一边界，平盘最便于计算；全微分项可先消去。\par\textbf{变式与易错点} 反向绕行时三个方法的结果都应变号。\end{example}

\section{曲线方向与未知函数积分}\label{ic:delta_fields}

\begin{examtip}本节用有向空间曲线与测试函数训练方向辨识，并通过分部公式消去未知边值。\end{examtip}

\subsection{空间交线参数与补线核对}\label{ic:delta_fields:spaceorientation}\begin{method}{空间交线参数与补线核对}\textbf{识别信号} 曲线由两个曲面相交且指定起终点及象限。\par \textbf{操作步骤} 先用平面消一个变量，参数化剩余椭圆；由象限选参数区间；补线时其方向必须使边界闭合。\par \textbf{疑难点与条件} 同一对端点可能对应两条弧，象限条件决定正确弧。\end{method}

\begin{example}{竞赛风格｜交线弧的方向与两个算法}{难度：$\star\star\star$}\label{ic:ex:delta_fields:spacearc}\textbf{题目} $L$为球$x^2+y^2+z^2=1$与平面$x+z=1$在第一卦限的弧，从$(1,0,0)$到$(0,0,1)$。求$\int_L(y\dd x+z\dd y+x\dd z)$。\par \textbf{思路} 椭圆参数后直接积分，并与补线Stokes核对。\par \textbf{详细解答} 取$x=\frac{1+\cos t}{2},z=\frac{1-\cos t}{2},y=\frac{\sin t}{\sqrt2}$，$0\le t\le\pi$。代入得到三项积分：$\int y\dd x=-\frac{\pi}{4\sqrt2}$，$\int z\dd y=-\frac{\pi}{4\sqrt2}$，$\int x\dd z=\frac{1}{2}$。总值$\frac{1}{2}-\frac{\pi}{2\sqrt2}$。补从终点到起点的线段$y=0,z=1-x$，其积分$-\int_0^1x\dd x=-\frac{1}{2}$；闭路的平面半椭圆实际是半径$\frac{1}{\sqrt2}$的半圆，面积$\frac{\pi}{4}$；相容法向为$\frac{1}{\sqrt2}(1,0,1)$，旋度$(-1,-1,-1)$与法向点乘为$-\sqrt2$，故闭路积分由Stokes为$-\frac{\pi}{2\sqrt2}$，减去补线同得结果。\par \textbf{验算与条件复核} 直接法使用$\int_0^\pi\sin^2t=\int_0^\pi\cos^2t=\frac{\pi}{2}$与$\int_0^\pi\sin t=2$，逐项验证符号。\par \textbf{疑难点与方法比较} 补线的方向是从终点回起点；取另一半椭圆会违反$y\ge0$。\end{example}

\subsection{由积分分部构造零边界测试函数}\label{ic:delta_fields:testfunction}\begin{method}{由积分分部构造零边界测试函数}\textbf{识别信号} 只知道$\Delta f$而不知道边界上的$f$或法向导数。\par \textbf{操作步骤} 找$g$使$\nabla g$等于目标积分中的向量权重，且$g$在边界为0；用$\int\nabla g\cdot\nabla f=-\int g\Delta f$消去未知边值。\par \textbf{疑难点与条件} 零边界的是测试函数$g$，不需要误加$f=0$的条件。\end{method}

\begin{theorem}{Green第一恒等式的推导}对有光滑边界的$\Omega$和足够光滑$f,g$，$\int_\Omega\nabla g\cdot\nabla f=\int_{\partial\Omega}g\,\partial_nf-\int_\Omega g\Delta f$。\end{theorem}\begin{proof}乘积法则$\operatorname{div}(g\nabla f)=\nabla g\cdot\nabla f+g\Delta f$，在$\Omega$积分再用Gauss即可。\end{proof}

\begin{example}{竞赛拓展｜不用解PDE也能求积分}{难度：$\star\star\star$}\label{ic:ex:delta_fields:laplacianmoment}\textbf{题目} $f\in C^2(\overline B)$，$B$是单位球，满足$\Delta f=\sqrt{x^2+y^2+z^2}$。求$\iiint_B(xf_x+yf_y+zf_z)\dd V$。\par \textbf{思路} 取$g=\frac{r^2-1}{2}$，它在球面为0且$\nabla g=(x,y,z)$。\par \textbf{详细解答} Green第一恒等式使边界项为0，目标$=-\frac12\iiint_B(r^2-1)r\dd V=2\pi\int_0^1(1-r^2)r^3\dd r=2\pi(\frac{1}{4}-\frac{1}{6})=\frac{\pi}{6}$。\par \textbf{验算与条件复核} 测试函数的梯度、边界值和球坐标体元$r^2$分别核对后，所有未知$f$消失。\par \textbf{疑难点与方法比较} 不能擅自假设$f$为径向函数；这里只用到给出的Laplacian。\end{example}

\section{图解：边界方向与法向}
\begin{figure}[htbp]\centering
\begin{tikzpicture}[scale=1.2,>=stealth]
\fill[Teal!8](0,0)circle(1.6);
\draw[Ink,thick](0,0)circle(1.6);
\draw[->,Teal,very thick](0:1.6)arc(0:80:1.6);
\draw[->,Teal,very thick](180:1.6)arc(180:260:1.6);
\node at(0,0){$\odot\ \bm n$};
\node[align=left,anchor=west]at(2.2,.3){法向指向读者\\正向边界逆时针\\法向反转，边界也反转};
\end{tikzpicture}
\caption{平面Green公式的正向与Stokes的右手规则一致。对单位圆及$\bm F=(-y,x,0)$，$\nabla\times\bm F=(0,0,2)$，正向积分$2\pi$；反向则为$-2\pi$。图中的点圈表示法向朝向读者。}\label{fig:review:orientation}
\end{figure}

\section{曲面度量的消因子、切向消项与逆向识别}\label{cf:section}
\begin{examtip}
曲线与曲面积分可在参数化之前先作几何化简：隐式法向量能消去曲面度量，锥面上的径向向量属于切空间，路径族把积分最值化成一元问题。最后从恰当条件反求未知函数，并检查穿孔区域中是否真的存在非零环流。
\end{examtip}

\subsection{隐式曲面的投影消因子}\label{cf:projection}
\begin{method}{识别法向余弦与面积元的乘积}
\textbf{适用条件} $F(x,y,z)=0$是正则曲面，且某一支可投影为单值图形，$F_z\ne0$。\par
\textbf{步骤} 用$z_x=-\frac{F_x}{F_z},z_y=-\frac{F_y}{F_z}$得到
\[
\dd S=\dfrac{|\nabla F|}{|F_z|}\dd x\dd y,
\qquad
\dfrac{|F_z|}{|\nabla F|}\dd S=\dd x\dd y.
\]
在选定的一支上先消去这对因子，再计算投影区域。\par
\textbf{疑难点} 绝对值对应第一类面积元；若整个曲面在同一投影点上有多层，必须逐支计算，不能只算一次投影。
\end{method}
\begin{example}{竞赛风格｜倾斜椭球上的度量抵消}{难度：$\star\star\star$}
\label{cf:ex:projection}
\textbf{题目} 设
$F=x^2+y^2+z^2-yz-1$，$\Sigma$为$F=0$上满足$2z-y\ge0$的一支。求
\[
I=\iint_\Sigma
\dfrac{(x+\sqrt3)|2z-y|}
{\sqrt{4+y^2+z^2-4yz}}\dd S.
\]
\textbf{思路} 分子中的$2z-y$是$F_z$，分母是曲面上$|\nabla F|$。\par
\textbf{详细解答} 有
\[
\begin{aligned}
|\nabla F|^2
&=4x^2+(2y-z)^2+(2z-y)^2\\
&=4(x^2+y^2+z^2-yz)+y^2+z^2-4yz\\
&=4+y^2+z^2-4yz.
\end{aligned}
\]
配方得$(z-\frac{y}{2})^2=1-x^2-\frac{3y^2}{4}$，故该支是
$z=\frac{y}{2}+\sqrt{1-x^2-\frac{3y^2}{4}}$，投影为
$D=\{x^2+\frac{3y^2}{4}\le1\}$。
在投影内部$F_z>0$，由消因子公式
\[
I=\iint_D(x+\sqrt3)\dd A
=\sqrt3\,\operatorname{Area}(D)
=\sqrt3\cdot\pi\cdot1\cdot\dfrac2{\sqrt3}=2\pi.
\]
投影边界对应$F_z=0$，是一条曲线；可先在内部截去边界邻域，再取极限，结果不受该零面积边界影响。\par
\textbf{方法点评} 无须显式展开复杂的$\sqrt{1+z_x^2+z_y^2}$；隐式法向量更快地暴露抵消结构。\par
\textbf{易错点} 若同时取上下两支，不能仍只积一次$D$；也不能将$F_z=0$误认成$\nabla F=0$。
\end{example}

\subsection{切向分量对通量没有贡献}\label{cf:tangent}
\begin{method}{先检查场是否沿曲面切向}
\textbf{适用条件} 计算通量，场的一部分在曲面上与法向量正交。\par
\textbf{步骤} 把场分解成切向项和剩余项，切向项的法向分量为零，可在积分前消去。\par
\textbf{疑难点} 消去依赖曲面几何，不要求这一部分场在整个体域散度为零；未知函数只有连续性时，直接投影比套散度定理更合适。
\end{method}
\begin{example}{考研风格｜锥面上未知函数自动消失}{难度：$\star\star\star$}
\label{cf:ex:tangent}
\textbf{题目} 设$h$为连续函数，$\Sigma$为锥面$z=\sqrt{x^2+y^2}$上$1\le x^2+y^2\le4$的下侧。求场
\[
\bm F=h(xy)(x,y,z)+(2x-y,\ 2y+x,\ z)
\]
穿过$\Sigma$的通量。\par
\textbf{思路} 锥面的径向向量$(x,y,z)$是切向，故含$h$的项不产生通量。\par
\textbf{详细解答} 令$r=\sqrt{x^2+y^2}$，下侧定向面积元为
$\bm n\dd S=(\frac{x}{r},\frac{y}{r},-1)\dd x\dd y$。
在曲面上$z=r$，有
\[
h(xy)(x,y,z)\cdot(\dfrac{x}{r},\dfrac{y}{r},-1)=h(xy)(r-z)=0.
\]
剩余部分点乘为
\[
\dfrac{x(2x-y)+y(2y+x)}r-z=2r-z=r.
\]
所以
\[
\iint_\Sigma\bm F\cdot\bm n\dd S
=\int_0^{2\pi}\int_1^2r^2\dd r\dd\theta
=\dfrac{14\pi}{3}.
\]
另一解释是对锥面的齐次方程$z-r=0$使用Euler恒等式，得到径向向量与法向量正交。\par
\textbf{方法点评} 几何正交使未知函数整体消失，不必求$h'$，也不必补面。\par
\textbf{易错点} 下侧面积元的第三分量为$-1$；反向取上侧，结果为$-\frac{14\pi}{3}$。
\end{example}

\subsection{路径族的积分最值}\label{cf:pathfamily}
\begin{method}{把路径参数代入后作一元优化}
\textbf{适用条件} 路径由一个或有限个实参数确定，起终点固定。\par
\textbf{步骤} 先保留路径方向，把曲线积分写成参数函数$I(a)$；再求参数可行域上的最值，必要时用Green补线核对。\par
\textbf{疑难点} 在某个路径族内最优不意味着在所有可能曲线中最优。
\end{method}
\begin{example}{竞赛风格｜正弦路径族中的最小做功}{难度：$\star\star\star$}
\label{cf:ex:pathfamily}
\textbf{题目} $L_a$为$y=a\sin x$、$0\le x\le\pi$，方向由$(0,0)$到$(\pi,0)$，$a>0$。求
\[
I(a)=\int_{L_a}(1+y^3)\dd x+(2x+y)\dd y
\]
在此路径族中的最小值及对应路径。\par
\textbf{思路} 代入$\dd y=a\cos x\dd x$，再对$a$优化。\par
\textbf{详细解答} 直接计算
\[
\begin{aligned}
I(a)
&=\int_0^\pi
[1+a^3\sin^3x+2ax\cos x+a^2\sin x\cos x]\dd x\\
&=\pi+\dfrac43a^3-4a.
\end{aligned}
\]
这里$\int_0^\pi\sin^3x\dd x=\frac{4}{3}$，
$\int_0^\pi x\cos x\dd x=[x\sin x+\cos x]_0^\pi=-2$，
最后一项积分为零。
因此$I'(a)=4(a^2-1)$，在$(0,1)$为负、$(1,\infty)$为正，唯一最优为$a=1$，
$I_{\min}=\pi-\frac{8}{3}$。

第二法补$x$轴上由$(\pi,0)$回$(0,0)$的线段，其积分为$-\pi$。
闭路为顺时针，Green给
\[
I(a)-\pi
=-\int_0^\pi\int_0^{a\sin x}(2-3y^2)\dd y\dd x
=-4a+\dfrac43a^3,
\]
与直接法相同。\par
\textbf{方法点评} 直接法便于显式求值，补线法解释了$-4a$与$a^3$项来自区域中的低阶矩。\par
\textbf{易错点} 上方弧从左向右行进、下方线段从右向左返回，整体是顺时针，Green公式要加负号。
\end{example}

\subsection{恰当条件反求未知场与全局势}\label{cf:inverse}
\begin{method}{用独立变量的系数比较反求未知函数}
\textbf{适用条件} 一个分量只含未知的一元函数，其余变量仍可自由变化。\par
\textbf{步骤} 先由恰当条件$P_y=Q_x$建立恒等式，再比较独立变量的系数确定未知函数；最后另查全局闭路与势函数。\par
\textbf{疑难点} 穿孔区域有洞，只表示无旋不足以保证有势，不表示任何无旋场都一定具有非零环流。
\end{method}
\begin{example}{竞赛风格｜有孔区域中仍可存在全局势}{难度：$\star\star\star$}
\label{cf:ex:inverse}
\textbf{题目} $\phi\in C^1(\R)$，在右半平面$x>0$上，一形式
\[
\omega=\dfrac{\phi(y)\dd x+2xy\dd y}{2x^2+y^4}
\]
是某函数的全微分。求$\phi$，并计算相应$\omega$沿穿孔平面$\R^2\setminus\{0\}$中任意分段光滑闭路的积分。\par
\textbf{思路} 右半平面的恰当条件确定$\phi$；计算闭路时还要研究角函数是否能全局连续选支。\par
\textbf{详细解答} 记$D=2x^2+y^4$。由$P_y=Q_x$得
\[
\phi'(y)(2x^2+y^4)-4y^3\phi(y)=2y^5-4x^2y.
\]
对所有$x>0$比较$x^2$系数，得到$\phi'=-2y$，故$\phi=-y^2+C$；
代回常数项得$4Cy^3=0$对所有$y$成立，所以$C=0$。
于是
\[
\omega=\dfrac{-y^2\dd x+2xy\dd y}{2x^2+y^4}
=\dfrac1{\sqrt2}\dd\theta,\qquad
\theta=\arg(\sqrt2x+\mathrm{i}y^2).
\]
此处取上半平面的辐角$0\le\theta\le\pi$。
当$y\ne0$时，可明确写成
$\theta=\frac{\pi}{2}-\arctan(\frac{\sqrt2x}{y^2})$。
当$y=0,x>0$时置$\theta=0$，当$y=0,x<0$时置$\theta=\pi$。
在前者附近$\theta=\arctan(\frac{y^2}{\sqrt2x})$，
后者附近$\theta=\pi+\arctan(\frac{y^2}{\sqrt2x})$，故这些表达拼成穿孔平面上的光滑函数。
直接求导有
\[
\theta_x=-\dfrac{\sqrt2y^2}{D},\qquad
\theta_y=\dfrac{2\sqrt2xy}{D}.
\]
因此$\frac{\theta}{\sqrt2}$是全局单值势，任意闭路积分均为$0$。\par
\textbf{方法点评} 映射$(x,y)\mapsto(\sqrt2x,y^2)$的像位于闭上半平面，绕行不会使所选辐角增加$2\pi$；这与通常极角场不同。\par
\textbf{易错点} 不能仅因分母在原点为零就断言闭路积分非零；也不能把单个$\arctan(\frac{y^2}{\sqrt2x})$分支不加修正地用于整个穿孔平面。
\end{example}
\begin{figure}[htbp]\centering
\begin{tikzpicture}[>=stealth]
\begin{scope}[shift={(1.8,1.5)}]
\draw[->,gray](-1.5,0)--(1.65,0)node[right]{$x$};
\draw[->,gray](0,-1.35)--(0,1.6)node[above]{$y$};
\draw[Ink,thick](0,0)circle(1.1);
\draw[->,Teal,very thick](1.1,0)arc[start angle=0,end angle=115,radius=1.1cm];
\draw[->,Gold,very thick](-1.1,0)arc[start angle=180,end angle=295,radius=1.1cm];
\fill[Ink](0,0)circle(1.2pt);
\node[below left]at(0,0){$O$};
\node[below,align=center,font=\small]at(0,-1.6){单位圆绕原点一周\\上、下半周依次经过};
\end{scope}
\draw[->,Ink,thick](4.1,1.55)--(6.3,1.55);
\node[above,align=center,font=\small]at(5.2,1.6){$u=\sqrt2x$\\$v=y^2$};
\begin{scope}[shift={(9,0.6)},x=1.35cm,y=1.7cm]
\draw[->,gray](-1.7,0)--(1.8,0)node[right]{$u$};
\draw[->,gray](0,-.2)--(0,1.45)node[above]{$v$};
\draw[Teal,thick,domain=-1.4142:1.4142,samples=65]plot(\x,{1-\x*\x/2});
\draw[->,Teal,very thick,domain=.95:.5,samples=15]plot(\x,{1-\x*\x/2});
\draw[->,Gold,very thick,domain=-.95:-.5,samples=15]plot(\x,{1-\x*\x/2});
\fill[Ink](0,0)circle(1pt);
\node[below]at(1.4142,0){$\sqrt2$};
\node[below]at(-1.4142,0){$-\sqrt2$};
\node[below left]at(0,0){$O$};
\node[above,font=\small]at(0,1.08){$v=1-\frac{u^2}{2}$};
\node[below,align=center,font=\small]at(0,-.48){映射后沿同一弧线往返\\不绕原点，辐角净变化为0};
\end{scope}
\end{tikzpicture}
\caption{原圆周的两半在$(u,v)$平面上重合，方向相反。图形解释了本例为何与通常极角场不同；任意闭路积分为零的严格依据仍是正文构造出的全局单值势。}\label{cf:fig:argumentmap}
\end{figure}

\section{闭路变形与奇点：先挖洞，再核对旋度}\label{r2:green}
设 $C$ 为正向简单闭曲线，区域内有孤立奇点 $a$。取围绕 $a$ 的小圆 $C_\varepsilon$，按逆时针方向记其积分。在去掉圆盘的区域上，若 $P,Q\in C^1$，Green 公式是
\[
 \oint_C P\dd x+Q\dd y-\oint_{C_\varepsilon}P\dd x+Q\dd y
 =\iint_{D\setminus B_\varepsilon(a)}(Q_x-P_y)\dd A.
\]
内边界在区域正向定向下是顺时针，因此当小圆单独按逆时针计时，上式有减号。只有去孔区域中 $Q_x-P_y=0$ 时，才能立即说大闭路积分等于小圆积分；有多个奇点时，等于各小圆积分的和。若旋度不为零，必须保留面积积分项。

\begin{example}{原点角度场：大圈与小圈为何同值}{数一方法训练}
取 $P=-y/(x^2+y^2),Q=x/(x^2+y^2)$。在原点外有
\[
 Q_x=P_y=\frac{y^2-x^2}{(x^2+y^2)^2}.
\]
因此对围住原点的正向简单闭曲线，可缩到逆时针小圆。令 $x=\varepsilon\cos t,y=\varepsilon\sin t$，被积微分化为 $\dd t$，故积分为 $2\pi$，与半径无关。若曲线内部不含原点且区域满足公式条件，则积分为零；方向反转时符号改变。
\end{example}

\begin{warning}
不能把奇点处不连续的场直接套到未挖洞的区域上，得出“旋度为零所以积分为零”。也不能认为任意奇点都可忽略：对大小为 $O(\varepsilon^{-p})$ 的场，小圆长为 $O(\varepsilon)$，只能估计圆周积分为 $O(\varepsilon^{1-p})$；$p=1$ 时并不保证极限为零。上述角度场就是非零例子。
\end{warning}

\paragraph{非零旋度例题与答案}
把上例 $Q$ 改为 $x/(x^2+y^2)+x$，则去孔区域旋度为一。对单位圆，原积分应为 $2\pi+\pi=3\pi$，而不是只取小圆的 $2\pi$。其中新增项 $\oint x\dd y=\pi$，可由参数化直接核对。

\paragraph{空间中的 Stokes 变形}
若两条定向边界 $C_1,C_2$ 围成一片场在其邻域为 $C^1$ 的有向连接曲面 $S$，并且 $\partial S=C_1-C_2$，则
\[
 \oint_{C_1}\bm F\cdot\dd\bm r-\oint_{C_2}\bm F\cdot\dd\bm r
 =\iint_S(\nabla\times\bm F)\cdot\bm n\dd S.
\]
要推出两闭路积分相同，须让右边为零，且变形曲面不能穿越奇点或奇异线。上面的角度场延伸为三维场 $(-y/(x^2+y^2),x/(x^2+y^2),0)$ 后，奇异集是整条 $z$ 轴，不能把绕轴圆周穿过轴缩成一点。


\section*{章末分层习题}\addcontentsline{toc}{section}{章末分层习题}
\refstepcounter{exercise}\label{q:8:1}\par\noindent\textbf{习题\theexercise\quad 基础题} 半径2圆周求$\int_L1\dd s$。\quad\hyperref[a:8:1]{[答案]}\par
\refstepcounter{exercise}\label{q:8:2}\par\noindent\textbf{习题\theexercise\quad 基础题} 求$\bm F=(x,y,z)$的散度。\quad\hyperref[a:8:2]{[答案]}\par
\refstepcounter{exercise}\label{q:8:3}\par\noindent\textbf{习题\theexercise\quad 基础题} 求$\bm F=(-y,x,0)$的旋度。\quad\hyperref[a:8:3]{[答案]}\par
\refstepcounter{exercise}\label{q:8:4}\par\noindent\textbf{习题\theexercise\quad 考研风格} 单位圆顺时针求$\oint x\dd y$。\quad\hyperref[a:8:4]{[答案]}\par
\refstepcounter{exercise}\label{q:8:5}\par\noindent\textbf{习题\theexercise\quad 考研风格} 求$\int_L\dd(xy)$，起点$(1,2)$终点$(3,4)$。\quad\hyperref[a:8:5]{[答案]}\par
\refstepcounter{exercise}\label{q:8:6}\par\noindent\textbf{习题\theexercise\quad 考研风格} 平面$z=1$单位圆盘上侧，恒场$(0,0,2)$的通量。\quad\hyperref[a:8:6]{[答案]}\par
\refstepcounter{exercise}\label{q:8:7}\par\noindent\textbf{习题\theexercise\quad 竞赛拓展} 单位球外侧求$(x^3,y^3,z^3)$的通量。\quad\hyperref[a:8:7]{[答案]}\par
\refstepcounter{exercise}\label{q:8:8}\par\noindent\textbf{习题\theexercise\quad 竞赛拓展} 证明$\oint_L\nabla f\cdot\dd\bm r=0$，$L$闭合且$f\in C^1$。\quad\hyperref[a:8:8]{[答案]}\par
\refstepcounter{exercise}\label{q:8:9}\par\noindent\textbf{习题\theexercise\quad 竞赛拓展} 解释穿孔平面无旋场不一定有全局势。\quad\hyperref[a:8:9]{[答案]}\par
\section*{章末方法速查}\addcontentsline{toc}{section}{章末方法速查}
\par\noindent\begin{tabularx}{\textwidth}{p{23mm}XXX}\toprule 方法&适用条件&典型信号&易错点\\\midrule
参数法&明确参数&简单曲线面&方向与模混淆\\Green&平面闭曲线&面积与环流&内部奇点\\Gauss&闭曲面外侧&散度简单&开面未补\\Stokes&定向边界&空间环流&右手方向反转\\
\bottomrule\end{tabularx}
